\appendixsection{随机森林算法关键代码}\label{sec:appendix-code}

本附录选取随机森林算法实现中的部分核心代码片段，以说明训练样本构造与在线预测服务的具体实现方式。为避免附录内容过长，此处仅保留与论文正文分析直接相关的关键逻辑，省略可视化输出及部分辅助性代码。

\appendixsubsection{训练样本生成与评分计算}

\cref{lst:rf-generate} 给出了训练样本生成过程中最核心的特征构造与评分计算逻辑。该代码用于模拟生成拼车候选订单的时空特征、用户信用特征以及目标评分，为随机森林模型训练提供结构化样本数据。

\begin{lstlisting}[language=Python,caption={随机森林训练样本生成核心代码},label={lst:rf-generate}]
def generate_sample():
    # 1. 生成空间匹配特征
    start_match = round(random.uniform(0, 1), 3)
    end_match = round(random.uniform(0, 1), 3)

    # 2. 生成时间匹配特征
    time_overlap = round(random.uniform(0, 1), 3)
    time_center_diff = random.randint(0, 120)
    is_time_compatible = 1 if time_overlap > 0.2 else 0

    # 3. 生成用户侧特征
    user_credit = round(random.uniform(0.5, 1), 3)
    success_rate = round(random.uniform(0.3, 1), 3)
    late_rate = round(random.uniform(0, 0.5), 3)

    # 4. 按预设权重计算目标评分
    score = (
        0.2 * start_match +
        0.2 * end_match +
        0.25 * time_overlap +
        0.1 * (1 / (1 + time_center_diff)) +
        0.1 * is_time_compatible +
        0.1 * user_credit +
        0.1 * success_rate -
        0.05 * late_rate
    )

    # 5. 添加扰动并裁剪到 [0,1]
    noise = random.uniform(-0.05, 0.05)
    score = max(0, min(1, score + noise))

    # 6. 返回完整样本
    return [
        start_match, end_match, time_overlap, time_center_diff,
        is_time_compatible, user_credit, success_rate, late_rate,
        round(score, 3)
    ]
\end{lstlisting}

\appendixsubsection{模型训练与评估实现}

\cref{lst:rf-train} 展示了随机森林模型训练阶段的核心实现。该部分代码负责完成数据读取、训练集与测试集划分、模型参数设置、误差指标计算以及训练结果保存，是拼车匹配算法从样本数据到可部署模型的关键环节。

\begin{lstlisting}[language=Python,caption={随机森林模型训练核心代码},label={lst:rf-train}]
# 数据文件路径与模型保存路径
DATA_PATH = "data.csv"
MODEL_PATH = "rf_model.pkl"

# 训练阶段的关键参数
TEST_SIZE = 0.2
RANDOM_STATE = 42
N_ESTIMATORS = 150
MAX_DEPTH = 12

def predict_score(model, features):
    # 将单个样本转换为二维数组
    features_array = np.array(features).reshape(1, -1)
    pred = model.predict(features_array)[0]

    # 将预测结果裁剪到评分区间
    if pred < 0:
        pred = 0.0
    elif pred > 1:
        pred = 1.0

    return float(pred)

def main():
    # 1. 读取训练数据
    df = pd.read_csv(DATA_PATH)

    # 2. 提取特征和标签
    X = df.iloc[:, 0:8].values
    y = df.iloc[:, 8].values

    # 3. 划分训练集与测试集
    X_train, X_test, y_train, y_test = train_test_split(
        X, y, test_size=TEST_SIZE, random_state=RANDOM_STATE
    )

    # 4. 初始化随机森林模型
    model = RandomForestRegressor(
        n_estimators=N_ESTIMATORS,
        max_depth=MAX_DEPTH,
        random_state=RANDOM_STATE
    )

    # 5. 训练并预测
    model.fit(X_train, y_train)
    y_pred = model.predict(X_test)

    # 6. 计算误差指标
    mae = mean_absolute_error(y_test, y_pred)
    rmse = np.sqrt(mean_squared_error(y_test, y_pred))

    print("MAE:", mae)
    print("RMSE:", rmse)

    # 7. 保存模型文件，供在线预测使用
    joblib.dump(model, MODEL_PATH)
    print("模型保存为:", MODEL_PATH)
\end{lstlisting}

\appendixsubsection{规则评分与特征封装实现}

\cref{lst:rule-score-java} 给出了拼车服务中规则评分的核心实现。代码先计算起点匹配、终点匹配和时间重叠等特征，再补入用户信用相关指标，最后按设定权重生成规则分数，并封装为 \texttt{RuleScore} 对象，为模型重排阶段提供输入。

\begin{lstlisting}[language=Java,caption={拼车规则评分与特征封装核心代码},label={lst:rule-score-java}]
@Data
@AllArgsConstructor
public class RuleScore {
    private CarpoolOrder order;
    private double score;
    private double startMatch;
    private double endMatch;
    private double timeOverlap;
    private double timeCenterDiff;
    private double isTimeCompatible;
    private double userCredit;
    private double successRate;
    private double lateRate;

    // 按固定顺序返回模型输入特征
    public double[] features() {
        return new double[]{
                startMatch,
                endMatch,
                timeOverlap,
                timeCenterDiff,
                isTimeCompatible,
                userCredit,
                successRate,
                lateRate
        };
    }
}

private RuleScore scoreOrder(CarpoolOrder order, String startLocation, String endLocation,
                             LocalTime startTime, LocalTime endTime) {
    // 1. 计算空间和时间匹配特征
    double startMatch = LocationMatcher.match(startLocation, order.getStartLocation());
    double endMatch = LocationMatcher.match(endLocation, order.getEndLocation());
    double timeOverlap = TimeMatchCalculator.overlapRatio(
            startTime, endTime, order.getStartTime(), order.getEndTime());
    double timeCenterDiff = TimeMatchCalculator.centerDiffMinutes(
            startTime, endTime, order.getStartTime(), order.getEndTime());
    double isTimeCompatible = timeOverlap > 0 ? 1.0 : 0.0;

    // 2. 获取发布者信用统计信息
    UserProfile profile = userClient.getUser(order.getUserId());
    double userCredit = defaultNumber(profile.getUserCredit(), 1.0);
    double successRate = defaultNumber(profile.getSuccessRate(), 1.0);
    double lateRate = defaultNumber(profile.getLateRate(), 0.0);

    // 3. 按预设权重生成规则分数
    double score = 0.2 * startMatch
            + 0.2 * endMatch
            + 0.25 * timeOverlap
            + 0.1 * (1 / (1 + timeCenterDiff))
            + 0.1 * isTimeCompatible
            + 0.1 * userCredit
            + 0.1 * successRate
            - 0.05 * lateRate;

    // 4. 返回封装后的候选评分对象
    return new RuleScore(
            order,
            clamp(score),
            startMatch,
            endMatch,
            timeOverlap,
            timeCenterDiff,
            isTimeCompatible,
            userCredit,
            successRate,
            lateRate
    );
}
\end{lstlisting}

\appendixsubsection{拼车反馈提交实现}

\cref{lst:feedback-java} 展示了拼车服务中反馈提交的关键处理逻辑。该实现负责校验评价资格、过滤非法评价对象、防止重复提交，并在保存反馈记录后同步更新用户拼车统计，是信用约束闭环中的重要环节。

\begin{lstlisting}[language=Java,caption={拼车反馈提交核心代码},label={lst:feedback-java}]
@Transactional
@SuppressWarnings("unchecked")
public Map<String, Object> submitFeedback(Long orderId, Long reviewerUserId,
                                          List<Map<String, Object>> feedbackItems) {
    // 1. 先确认订单当前是否允许反馈
    CarpoolOrder order = getOrder(orderId);
    if (!isFeedbackAvailable(order)) {
        throw new IllegalArgumentException("Feedback is not available for this order");
    }

    // 2. 计算当前用户可评价的目标集合
    List<Long> allowedTargets = feedbackTargetIds(order, reviewerUserId);
    if (allowedTargets.isEmpty()) {
        throw new IllegalArgumentException("No feedback target for this user");
    }
    if (feedbackItems == null || feedbackItems.isEmpty()) {
        feedbackItems = defaultFeedbackItems(allowedTargets);
    }

    int savedCount = 0;
    for (Map<String, Object> item : feedbackItems) {
        Long targetUserId = longValue(item.get("targetUserId"), null);

        // 3. 校验评价目标是否合法
        if (targetUserId == null || !allowedTargets.contains(targetUserId)) {
            throw new IllegalArgumentException("Invalid feedback target");
        }

        // 4. 同一订单中避免重复评价
        if (feedbackRepository.existsByOrderIdAndReviewerUserIdAndTargetUserId(
                orderId, reviewerUserId, targetUserId)) {
            continue;
        }

        double score = doubleValue(item.get("score"), 1.0);
        if (score < 0 || score > 1) {
            throw new IllegalArgumentException("score must be between 0 and 1");
        }

        boolean late = booleanValue(item.get("late"), false);
        boolean completed = booleanValue(item.get("completed"), true);

        // 5. 构造并保存反馈记录
        CarpoolFeedback feedback = new CarpoolFeedback();
        feedback.setOrderId(orderId);
        feedback.setReviewerUserId(reviewerUserId);
        feedback.setTargetUserId(targetUserId);
        feedback.setScore(score);
        feedback.setLate(late);
        feedback.setCompleted(completed);
        feedback.setCreatedAt(LocalDateTime.now());
        feedbackRepository.save(feedback);

        // 6. 同步更新被评价用户的拼车统计
        userClient.updateCarpoolStats(targetUserId, score, late, completed);
        savedCount++;
    }

    return Map.of("submitted", true, "savedCount", savedCount, "orderId", orderId);
}
\end{lstlisting}

\appendixsubsection{在线预测接口实现}

\cref{lst:rf-api} 展示了模型服务中用于输入校验与在线预测的关键实现。该部分代码基于 FastAPI 提供 \texttt{/predict} 接口，在接收后端传入的特征向量后完成数据合法性检查、模型调用和预测结果裁剪，是拼车匹配算法在线执行流程中的关键支撑部分。

\begin{lstlisting}[language=Python,caption={随机森林在线预测接口核心代码},label={lst:rf-api}]
def validate_features(features):
    # 1. 判空检查
    if features is None:
        return False, "features 不能为空"
    # 2. 类型检查
    if not isinstance(features, list):
        return False, "features 必须是数组"
    # 3. 长度检查
    if len(features) != 8:
        return False, "features 长度必须为 8"
    # 4. 数值检查
    for f in features:
        if not isinstance(f, (int, float)):
            return False, "features 中必须全部为数字"
    return True, None

@app.post("/predict")
async def predict(request: Request):
    # 读取请求体中的特征向量
    data = await request.json()
    features = data.get("features")

    # 输入校验失败则直接返回错误信息
    is_valid, error_msg = validate_features(features)
    if not is_valid:
        return JSONResponse(status_code=400, content={"error": error_msg})

    # 调用模型进行预测
    input_data = [features]
    result = model.predict(input_data)
    score = float(result[0])

    # 将预测值裁剪到 [0,1]
    if score < 0:
        score = 0.0
    elif score > 1:
        score = 1.0

    return {"score": score}
\end{lstlisting}
